\begin{lemma}[The fractional spectral gap]
For every $v \in L^2(\gamma_n)$,
\[
  \|\Pi_0 v\|_2^2 \leq \mathcal{E}(v)
\]
as an inequality in $[0,+\infty]$, where $\Pi_0$ is as in \noderef{FracPow} and
$\mathcal{E}$ as in \noderef{Energy}.  Equivalently, by
\noderef{ProjZeroCentering},
$\|v - \mathfrak{m}(v)\mathbf{1}\|_2^2 \leq \mathcal{E}(v)$ with
$\mathfrak{m}(v) = \int v \, d\gamma_n$; in particular the inequality is vacuously
true when $\mathcal{E}(v) = +\infty$, and it is the spectral gap estimate of the
reference paper when $v$ lies in the form domain $\mathcal{V}$.
\end{lemma}
\begin{proof}
\emph{Step 1: the left-hand side as a coefficient sum.}
By \noderef{FracPowCoeff} and \noderef{FracPow} with $s = 0$, the element
$\Pi_0 v$ has Hermite coefficients $0$ at $\alpha = 0$ and $v_\alpha$ at
$\alpha \neq 0$ (the domain condition of \noderef{FracPowDomain} being satisfied,
as the multiplied family is dominated by the square-summable family
$(|v_\alpha|^2)_\alpha$).  Parseval's identity from \noderef{HermiteBasis}
therefore gives
\[
  \|\Pi_0 v\|_2^2 = \sum_{\alpha \neq 0} |v_\alpha|^2 .
\]

\emph{Step 2: a termwise comparison.}
Let $\alpha \neq 0$.  By \noderef{MultiIndexSum}, $|\alpha| = \sum_j \alpha_j$ is
a nonnegative integer which vanishes only for $\alpha = 0$; hence
$|\alpha| \geq 1$, and so $|\alpha|^{1/2} \geq 1$.  Multiplying the nonnegative
number $|v_\alpha|^2$ by this gives
\[
  |v_\alpha|^2 \leq |\alpha|^{1/2}\,|v_\alpha|^2 .
\]

\emph{Step 3: summation.}
All terms on both sides are nonnegative, so summing the Step-2 inequalities over
$\alpha \in \mathbb{N}_0^n \setminus\{0\}$ is legitimate in $[0,+\infty]$ and
preserves the inequality:
\[
  \sum_{\alpha \neq 0} |v_\alpha|^2
    \;\leq\; \sum_{\alpha \neq 0} |\alpha|^{1/2}|v_\alpha|^2 .
\]
By \noderef{Energy} the right-hand side is exactly $\mathcal{E}(v)$, interpreted in
$[0,+\infty]$.  Combining with Step 1 gives
$\|\Pi_0 v\|_2^2 \leq \mathcal{E}(v)$.

Finally, by \noderef{ProjZeroCentering}, $\Pi_0 v = v - \mathfrak{m}(v)h_0$ with
$h_0 = \mathbf{1}$, so the left-hand side also equals
$\|v - \mathfrak{m}(v)\mathbf{1}\|_2^2$, giving the centred form.
\end{proof}
